% Source: https://artofproblemsolving.com/community/c3h2984662_natural_divisors
% Stable identity: aops-topic-2984662-q01
Write $n=13k$ and let $\tau$ denote the number of divisors; we need $\tau(13k)=k$.

\emph{Case $13\nmid k$.} Then $\tau(13k)=2\tau(k)$, so $\tau(k)=\dfrac k2$. Pairing each divisor $d<\sqrt k$ with $\dfrac kd$ shows $\tau(k)\le2\sqrt k$, so $\dfrac k2\le2\sqrt k$ and $k\le16$. Checking even $k\le16$: $\tau(k)=\dfrac k2$ holds for $k=8$ ($\tau=4$) and $k=12$ ($\tau=6$) only.

\emph{Case $k=13^am$ with $a\ge1$, $13\nmid m$.} Then $\tau(13k)=(a+2)\tau(m)\le(a+2)\cdot2\sqrt m$, which must equal $13^am\ge13^a\sqrt m$. So $2(a+2)\ge13^a$, impossible for $a\ge1$.

Hence $n=13\cdot8=104$ or $n=13\cdot12=156$: $\boxed{n\in\{104,156\}}$.
